\chapter{Following calls and saved state}
\MARSChapterHeading{Calls and saved state}
\label{revision:lab-functions}

This lab accompanies Chapter~5 of the main book. Use its
\texttt{reverse-rows.asm} worked program before attempting the separate
\texttt{reverse-and-invert.asm} exercise. The files are in
\texttt{revision/labs/files}. These activities use ordinary MARS Run and Debug;
this fifth revised unit is not yet a native interactive lesson.

Your evidence should connect a value to the call over which it must remain
available. A final memory dump can show two correct results, but it cannot by
itself explain how the caller's result pointer survived an inner function's use
of the same register.

\section{Separate a result from a continuation}

With the Default layout and MIPS64 Release 6 selected, break at
\texttt{main\_call}. Record \texttt{\$a0}, \texttt{\$a1}, \texttt{\$s0} and
\texttt{\$sp}. The argument words are \texttt{0x16} and \texttt{0x81}; the saved
register holds an address, not a third row value. The stack pointer has already
been aligned by the driver's two opening instructions.

Predict where the call will put \texttt{\$ra}. Compare the addresses in the
Disassembly view rather than counting source lines: an assembler convenience
such as \texttt{dli} may represent several machine instructions. The label
\texttt{after\_pair} names the immediate continuation of this compact call.

Use \emph{Step Into} to enter \texttt{reverse\_pair}. Check that execution reaches
its first instruction before the stores in \texttt{main}. \emph{Step Over} can
run a call and pause at its continuation; it is useful when testing an already
understood function, but it hides the frame construction we want to inspect here.

\section{Account for every saved byte}

Continue to \texttt{pair\_frame\_ready}, before the next \texttt{move}. Record the
new pointer and subtract it from the caller's pointer. The difference should be
32 bytes. Inspect the eight-byte slots at offsets 8, 16 and 24. Compare them with
the incoming \texttt{\$s1}, \texttt{\$s0} and \texttt{\$ra}. Interpret each slot
as a complete 64-bit value, not just the first displayed word.

In this layout, the aligned caller pointer is \texttt{0x7FFFEFF0}; the frame base
is \texttt{0x7FFFEFD0}. The saved return address begins at
\texttt{0x7FFFEFE8}. A slot beginning there occupies eight consecutive bytes and
ends at \texttt{0x7FFFEFEF}. The caller's first byte at
\texttt{0x7FFFEFF0} is outside the frame.

These addresses describe this run, not constants to put into the function.
Its offset-based accesses also work when a caller supplies a different valid,
aligned stack pointer. Explain why loading a saved slot after restoring the
pointer would address a different location.

\section{Follow two calls through the same helper}

At the first visit to \texttt{reverse\_byte}, record \texttt{\$ra}. It now names
\texttt{pair\_first\_ready}, while the outer return address remains in memory.
The helper does not reserve another frame because its temporary state fits in
registers and it calls no other function. Predict the eight iterations for the
first byte, then compare the returned \texttt{0x68} with your trace.

At \texttt{pair\_second\_call}, find the retained second input in both the saved
register and the prepared argument. Find the retained first result in
\texttt{\$s1}. Step into the same helper again. Its instruction addresses are
unchanged, but its input and return continuation differ from the first invocation.

At \texttt{pair\_second\_ready}, the second result is in \texttt{\$v0}. Predict
the two moves that place it in \texttt{\$v1} and restore the first result to
\texttt{\$v0}. Then inspect \texttt{pair\_return}: the caller's registers and
pointer must be restored before the return executes.

Continue to \texttt{after\_pair} and inspect the still-pending result stores.
Only then run to exit and inspect both result words and the guard. This separates
three claims: the function returned its results, it preserved the caller's state,
and the caller stored the results at its intended addresses.

\section{Investigate one lost value at a time}

Work on a copy when introducing a fault. First omit the restore of
\texttt{\$s0}. Before returning, identify exactly which promise has already
failed. The next store in \texttt{main} may fault or reach an unintended location,
but the saved-register mismatch is earlier evidence of the cause.

Next restore that instruction and instead omit saving the outer return address.
Both inner reversals may finish. The final return then lacks the continuation of
\texttt{main}; it is not enough that the loop calculated the right byte. Stop a
run that fails to reach the expected continuation and inspect the link's last
writer.

For a third copy, save a register with a word store and restore it with a word
load. A value with zero upper bits might hide that error. Try an incoming saved
value whose upper half is nonzero, using the debugger's register editor before
entering the function. Record why the preservation contract covers all 64 bits,
even though the row data uses only eight.

\section{Write a related function independently}

The independent file supplies a caller, a visible continuation and the same
\texttt{reverse\_byte} helper. Replace the failing \texttt{break} body with
\texttt{reverse\_and\_invert}. It must use the helper and return a result in
\texttt{\$v0} from zero through 255. It must return to its incoming continuation,
restore \texttt{\$sp}, preserve the convention's saved registers and leave
caller-owned stack bytes unchanged. It must not print or perform device I/O.

Begin with a call-based design: decide where to retain the incoming return
address, reserve an aligned frame if needed, call the helper, invert the returned
low byte, and restore what you saved. Do not copy a 32-byte frame automatically;
list the values your new routine actually needs to preserve.

For input \texttt{0x16}, the caller should retain \texttt{0x97} in
\texttt{\$s0} after the function returns. This assignment of \texttt{\$s0} is
performed by the supplied caller, after the callee's preservation obligation ends.
The exit service reuses \texttt{\$v0}; inspect the return boundary or the caller's
retained result rather than interpreting the final service selector as the answer.

Run the public checks with JDK~25 and the MARS CLI:

\begin{verbatim}
java -jar /path/to/cli.jar grade \
  revision/labs/checks/reverse-and-invert.kts \
  /path/to/your/reverse-and-invert.asm
\end{verbatim}

The public check exercises every low-byte input, adds nonzero upper argument bits,
and makes repeated calls with full-width saved-register values. It verifies the
result, execution of the provided helper, the restored pointer, saved registers,
caller-owned stack bytes and silent
execution. These are formative checks you can read. A pass does not replace your
explanation of the call boundaries.

An equivalent implementation may use different temporaries and frame slots. As an
extension, consider inverting the input byte first and transferring directly to
\texttt{reverse\_byte}, with the original return address still in place. Prove
that reversing inverted bits gives the same result as inverting reversed bits.
Explain why that final transfer would need no new return continuation. Test it
against the same contract rather than assuming that a shorter routine is correct.
